% 增加相应文献
% 1) 问题场景和挑战：真实 RAG 系统面对的是一个持续到来的 query stream，其中既有 long-tail 问题，也有当前 KB 根本没有证据的 out-of-support / unanswerable 问题；同时 KB 还受预算约束，不能无限增长。

% 2) 现有工作为什么解决不了：
% 现有动态 RAG 只在局部优化：有的只会增量加索引，有 的只会删/修 corpus，有的只会保留最近或高频记忆；但它们没有把 out-of-support queries 当成“KB 该怎么补”的信号。
% 它们分别优化了索引更新、知识编辑和历史缓存，但没有闭环解决：当当前 KB 无法覆盖新的 long-tail / out-of-support query 时，系统应如何从外部环境中主动补知识并重构 KB。
%    (a) 第一类：index-level incremental updates这类方法解决的是“怎么把新内容更快并进索引”，不是“怎么根据用户需求决定该补什么”。LightRAG 的重点是图索引、双层检索和增量更新；HippoRAG 2 的重点是长时记忆式的非参数 continual learning。它们能提升检索或记忆能力，但核心信号仍然是新增数据如何进索引，而不是query 分布如何驱动 KB 变更。
%    (b) 第二类：ingestion-time editing / pruning RECIPE、Zero-RAG 这类方法解决的是“怎么编辑/压缩已有知识库”，本质上是 supply-centric：要么修正错误，要么去掉冗余，要么把知识压成更适合推理的形式。它们关注的是 corpus 质量或冗余度，而不是“当前用户到底缺什么知识”。所以它们能让 KB 更干净，但不能自动把 out-of-support query 转成补知识动作。
%    (c) 第三类：temporal memory / caching / streaming ComRAG、RAGCache、StreamingRAG 解决的是“如何在时间维度上更快、更稳地保留已有记忆/缓存”。ComRAG 用 centroid-based memory 管理动态 QA 历史；RAGCache 用 PGDSF 做多级缓存；StreamingRAG 用 mini-batch clustering 和 heavy-hitter filter 维持紧凑原型。它们都很强，但优化对象是已有内容的保留与访问效率，不是从文档池里选择新的知识来修补 KB 覆盖缺口。

% =========================================================
% Motivation 实验需求文档（中文注释版）
% =========================================================
%
% 实验目标：
% ---------
% 本实验不是最终性能对比，而是一个“诊断实验（diagnostic experiment）”。
%
% 目标是证明：
% 现有动态 RAG 方法无法同时做到：
%   (1) 在长尾 / out-of-support 查询上保持覆盖能力
%   (2) 在固定 KB 预算下控制更新成本
%
% 更核心的是验证：
%   现有方法没有利用 query 侧的“知识不足信号”
%   来决定 是否更新 KB 以及 更新什么内容。
%
%
% =========================================================
% 1. 场景设定（Streaming Query Setting）
% =========================================================
%
% 构造一个 query stream：
%   q_1, q_2, ..., q_T 按时间顺序到达
%
% 每个 query 需要有两个标注：
%   - support label：
%       当前 KB 是否包含支持该 query 的证据
%   - tail label：
%       是否属于长尾（低频）query
%
% 同时对 KB 加预算约束 B：
%   - 最大文档数 / token 数 / index size
%
% KB 会随着时间动态更新（由不同策略决定）
%
%
% =========================================================
% 2. 核心问题（Motivation 核心）
% =========================================================
%
% 本实验关注的不是：
%   “query 和 KB 的相似度是否足够高”
%
% 而是：
%   当 query 暴露出“KB 无法支持”时：
%       系统是否利用这个信号去更新 KB？
%
% 关键缺口：
%   现有方法大多是：
%       - 文档驱动（document-driven）
%       - 或缓存驱动（cache-driven）
%
%   而不是：
%       - query 驱动（query-driven KB evolution）
%
%
% =========================================================
% 3. 现有更新范式总结（Representative Paradigms）
% =========================================================
%
% (A) 索引级增量更新（Index-level Incremental Updates）
%     - 触发方式：新文档到来就加入 KB
%     - 特点：append-only / 增量构建索引
%
%     局限：
%       * 内存随时间持续增长
%       * 不考虑 query 分布
%       * 不知道哪些知识“有用”
%
%
% (B) 语料编辑 / 剪枝（Corpus Editing / Pruning）
%     - 触发方式：文档到来时进行编辑、去重、纠错
%
%     局限：
%       * supply-driven（由数据驱动）
%       * 优化的是 KB 质量，而不是 query 覆盖
%       * 无法回答：应该补充哪些缺失知识
%
%
% (C) 时序记忆 / 缓存（Temporal Memory / Caching）
%     - 触发方式：基于访问频率 / 最近使用 / streaming 规则
%
%     局限：
%       * 强 recency bias（偏向最近/高频）
%       * 无法建模长期兴趣变化
%       * 无法主动引入缺失知识
%
%
% 可选：
% (D) 拒答策略（Reject-only）
%     - 低置信度时拒答
%
%     局限：
%       * 提高安全性，但不会提升 KB 能力
%
%
% =========================================================
% 4. Baseline 设计
% =========================================================
%
% Motivation 实验中只需选择“机制代表”：
%
%   1) Append-only（代表增量索引）
%   2) Editing / Pruning（代表语料编辑）
%   3) Cache / Recency（代表时序记忆）
%
% 可选：
%   4) Reject-only（作为对照）
%
% 注意：
%   这里不要求复现完整 SOTA 系统，
%   重点是对比“更新机制”。
%
%
% =========================================================
% 5. 指标设计（最小化指标集合）
% =========================================================
%
% 只使用 2~3 个核心指标：
%
% (1) Coverage@B
%     在预算 B 下，KB 能支持的 query 比例
%
% (2) OOS False Answer Rate
%     在 KB 无证据的 query 上：
%     系统仍然生成答案（而不是拒答）的比例
%
% (3) Update Cost（选一个即可）
%     - KB size（推荐）
%     - 或 update latency
%
% 可选：
%   Tail Coverage（仅在强调长尾时使用）
%
%
% =========================================================
% 6. 预期现象（Expected Findings）
% =========================================================
%
% Append-only：
%   + Coverage 可提升
%   - 成本持续增长
%
% Editing / Pruning：
%   + KB 更干净
%   - 不会补充 query 所需的新知识
%
% Cache / Recency：
%   + 对高频/近期 query 有效
%   - 长尾覆盖差
%
% Reject-only：
%   + 幻觉低
%   - 覆盖无提升
%
%
% =========================================================
% 7. 图设计（Figure 建议）
% =========================================================
%
% 推荐一个三联图：
%
%   (A) Coverage@B 随时间变化
%   (B) OOS False Answer Rate
%   (C) KB size / cost 增长
%
% 或简化为：
%   x轴：KB size
%   y轴：Coverage@B
%
%
% =========================================================
% 8. 核心结论（一句话）
% =========================================================
%
% 现有动态 RAG 方法的更新是“被动触发”的：
%   - 来自新文档
%   - 或缓存/访问模式
%
% 而不是由 query 暴露的知识缺口驱动。
%
% 因此，它们无法在预算约束下
% 持续对齐用户需求分布。
%
% =========================================================


% 5) 算法：
%    基于上述 observation，把候选文档看成对多个 interest centers 的覆盖问题，
%    用子模目标选 ideal KB，再受 budget 约束做增量替换。
%    这里要强调：算法不是拍脑袋规则，而是和 observation 一一对应。
\section{Introduction}

Retrieval-augmented generation (RAG) has substantially improved large language models (LLMs) on knowledge-intensive tasks~\cite{lewis2020rag}. In real deployments, however, the central difficulty is not whether retrieval helps in general, but how a system should respond when its current evidence is insufficient. Recent work has established that these failures are common. UAEval4RAG~\cite{peng2025uaeval4rag} shows that unsupported queries must be evaluated explicitly rather than absorbed into aggregate accuracy. FreshLLMs~\cite{vu2023freshllms} and LiveSearchBench~\cite{zhou2025livesearchbench} show that temporally novel knowledge remains difficult even with retrieval.

A well-studied line of work responds to these failures by keeping the KB \emph{static} while improving what happens at query time: retrieval-time correction, gating, and on-demand external search~\cite{yan2024crag,asai2023selfrag,vu2023freshllms,kasai2024realtime}. These methods are effective for the current query, but they leave the persistent KB unchanged. In many production deployments this is not sufficient. Systems with bounded storage, constrained or disallowed external lookup, latency budgets, or governance requirements must maintain a persistent, budget-limited KB and improve it over time. We study this well-motivated but under-solved operational regime.

Prior work on persistent KB maintenance has investigated three directions. Index-level incremental methods such as LightRAG~\cite{guo2024lightrag} and HippoRAG~2~\cite{guti2025hipporag2} optimize how newly arriving content is integrated into graph memory or index structures. Data-centric editing and pruning methods such as RECIPE~\cite{chen2024recipe} and Zero-RAG~\cite{zero-rag2025} improve the accuracy or compactness of existing knowledge. Temporal memory and caching methods such as ComRAG~\cite{comrag2025}, RAGCache~\cite{jin2024ragcache}, and StreamingRAG~\cite{wu2025streamingrag} optimize retention and reuse efficiency of admitted content. These are genuine contributions, but they share a common gap: none of them uses recurring query-exposed support failures as the primary signal for deciding what to admit, retain, or evict under a fixed KB budget. As a result, they leave three measurable shortfalls unaddressed. First, they exhibit high \emph{wasted update rate}: admitted content frequently fails to contribute to future unsupported queries because admission decisions are driven by supply-side availability or recency rather than demand-side gaps. Second, they achieve low \emph{future support gain per unit budget}: because they do not model which missing knowledge recurs across future queries, their budget allocation is inefficient relative to the coverage gaps they are meant to close. Third, they deteriorate on \emph{tail query coverage} under demand drift: when the query distribution shifts toward low-frequency or out-of-distribution topics, supply-driven and recency-biased methods have no mechanism to realign the KB toward emerging support gaps.

This paper addresses the resulting decision problem directly. Our central hypothesis is that persistent KB maintenance should be driven by query-exposed support failures as a first-class signal for selective KB evolution. We cast budgeted KB maintenance as an alignment problem between recurring query demand and persistent evidence coverage, and ask: given a stream of queries where support failures are observable, which missing knowledge is worth admitting under budget, and which failures indicate an unanswerable request that no KB update should chase?
